\documentclass[11pt,a4paper]{article}
% Preâmbulo comum. Não duplicar pacotes nos módulos.
\usepackage[utf8]{inputenc}
\usepackage[T1]{fontenc}
\usepackage[brazil]{babel}
\usepackage{amsmath,amssymb,amsthm,mathtools}
\usepackage{geometry}
\geometry{a4paper,margin=2.3cm}
\usepackage{microtype}
\usepackage{enumitem}
\usepackage{graphicx}
\usepackage{xcolor}
\usepackage{booktabs}
\usepackage{hyperref}
\hypersetup{colorlinks=true,linkcolor=blue!40!black,urlcolor=blue!50!black}

\setcounter{tocdepth}{2}
\numberwithin{equation}{section}

\theoremstyle{plain}
\newtheorem{teorema}{Teorema}[section]
\newtheorem{proposicao}[teorema]{Proposição}
\newtheorem{lema}[teorema]{Lema}
\newtheorem{corolario}[teorema]{Corolário}

\theoremstyle{definition}
\newtheorem{definicao}[teorema]{Definição}
\newtheorem{exemplo}[teorema]{Exemplo}
\newtheorem{exercicio}{Exercício}[section]

\theoremstyle{remark}
\newtheorem{observacao}[teorema]{Observação}
\newtheorem{contraexemplo}[teorema]{Contraexemplo}
\newtheorem{intuicao}[teorema]{Intuição}

\newcommand{\R}{\mathbb{R}}
\newcommand{\N}{\mathbb{N}}
\newcommand{\Z}{\mathbb{Z}}
\newcommand{\Q}{\mathbb{Q}}
\DeclareRobustCommand{\eps}{\varepsilon}
\DeclareRobustCommand{\vet}[1]{\vec{#1}}
\DeclareRobustCommand{\dif}{\mathrm{d}}
\DeclareRobustCommand{\abs}[1]{\left\lvert #1 \right\rvert}
\DeclareRobustCommand{\norm}[1]{\left\lVert #1 \right\rVert}
\DeclareMathOperator{\grad}{grad}
\DeclareMathOperator{\Div}{div}
\DeclareMathOperator{\rot}{rot}

\newcommand{\base}[1]{\par\medskip\noindent\textbf{Base para relembrar.} #1\par\medskip}
\newcommand{\cuidado}[1]{\par\medskip\noindent\textbf{Quando a hipótese falha.} #1\par\medskip}

\newcommand{\figura}[2]{%
  \begin{figure}[htbp]
  \centering
  \IfFileExists{#1}{%
    \includegraphics[width=0.84\textwidth]{#1}%
  }{%
    \fbox{\parbox{0.8\textwidth}{\centering Figura ainda não gerada:\\ \texttt{#1}}}%
  }
  \caption{#2}
  \end{figure}%
}

\newenvironment{esqueleto}{%
  \par\medskip\noindent\textbf{Esqueleto da demonstração.}\par\nopagebreak
  \begin{enumerate}[leftmargin=1.4em]
}{\end{enumerate}\medskip}

\title{Gabarito --- Desigualdades, valor absoluto e frações parciais}
\author{}
\date{}
\begin{document}
\maketitle
\tableofcontents

A numeração segue a ordem dos exercícios do texto. Cada enunciado é repetido e a solução vem em seguida.

\section*{Exercício 1. [Básico --- módulo e intervalo]}

\textbf{Enunciado.} Determine o conjunto $\{x\in\R:\abs{3x-6}\le 9\}$ e escreva-o como um intervalo.

\textbf{Solução.}
O fator $3$ é positivo, e $\abs{3(x-2)}=\abs{3}\,\abs{x-2}=3\abs{x-2}$. A desigualdade $3\abs{x-2}\le 9$, multiplicada pelo inverso positivo $1/3$, equivale a $\abs{x-2}\le 3$. Para $\delta=3\ge 0$, a caracterização do módulo dá
\[
-3\le x-2\le 3.
\]
Somando $2$ aos três membros, $-1\le x\le 5$. O conjunto solução é o intervalo fechado $[-1,5]$.

A mesma conta sem fatorar: $\abs{3x-6}\le 9$ equivale a $-9\le 3x-6\le 9$. Somando $6$, $-3\le 3x\le 15$. Multiplicando pelo positivo $1/3$, $-1\le x\le 5$.

\section*{Exercício 2. [Básico --- inversão por número negativo]}

\textbf{Enunciado.} Prove que, se $a<b$, então $-b<-a$, e exiba a inversão produzida por um fator negativo.

\textbf{Solução.}
Se $a<b$, a compatibilidade com a soma, aplicada ao número $-a-b$, dá
\[
a+(-a-b)<b+(-a-b),
\]
isto é $-b<-a$.

Para a segunda parte, $1<2$. O fator $-3$ satisfaz $-3<0$, então a hipótese ``o multiplicador é positivo'' do axioma do produto foi retirada. O item correspondente da ordem diz que o sentido se inverte: $1\cdot(-3)>2\cdot(-3)$, isto é $-3>-6$.

\section*{Exercício 3. [Básico --- mínimo por quadrado]}

\textbf{Enunciado.} Escreva $x^2-8x+3$ na forma $(x-h)^2+k$ e determine o menor valor em $\R$.

\textbf{Solução.}
O coeficiente de $x$ é $-8$, então $h=-4$:
\[
x^2-8x+3=(x-4)^2-16+3=(x-4)^2-13.
\]
Para todo real $t$, $t^2\ge 0$, com igualdade se e somente se $t=0$. Tomando $t=x-4$, $(x-4)^2\ge 0$, logo $(x-4)^2-13\ge -13$, com igualdade se e somente se $x-4=0$, isto é $x=4$. O menor valor é $-13$, atingido em $x=4$.

\section*{Exercício 4. [Básico --- frações lineares distintas]}

\textbf{Enunciado.} Decomponha $\dfrac{5}{(x-2)(x+3)}$ em frações parciais.

\textbf{Solução.}
Procura-se $A,B\in\R$ com
\[
\frac{5}{(x-2)(x+3)}=\frac{A}{x-2}+\frac{B}{x+3},\qquad x\neq 2,\ x\neq -3.
\]
Multiplicando por $(x-2)(x+3)$,
\[
5=A(x+3)+B(x-2)=(A+B)x+(3A-2B).
\]
O polinômio da direita é constante, então
\[
\begin{cases}
A+B=0,\\
3A-2B=5.
\end{cases}
\]
Da primeira equação, $A=-B$. Substituindo na segunda, $3(-B)-2B=5$, isto é $-5B=5$, logo $B=-1$ e $A=1$. A divisão por $-5\neq 0$ é equivalência, e a solução é única. Conferência:
\[
(x+3)-(x-2)=x+3-x+2=5.
\]
Assim
\[
\frac{5}{(x-2)(x+3)}=\frac{1}{x-2}-\frac{1}{x+3}.
\]

\section*{Exercício 5. [Intermediário --- três termos]}

\textbf{Enunciado.} Prove que $\abs{x+y+z}\le\abs{x}+\abs{y}+\abs{z}$ para todos $x,y,z\in\R$.

\textbf{Solução.}
Aplique a desigualdade triangular a $x+y$ e a $z$:
\[
\abs{(x+y)+z}\le\abs{x+y}+\abs{z}.
\]
Aplique de novo a triangular, agora a $x$ e a $y$: $\abs{x+y}\le\abs{x}+\abs{y}$. Somando $\abs{z}$ aos dois lados, a ordem não estrita preserva o sentido,
\[
\abs{x+y}+\abs{z}\le\abs{x}+\abs{y}+\abs{z}.
\]
A transitividade de $\le$ reúne as duas estimativas:
\[
\abs{x+y+z}\le\abs{x}+\abs{y}+\abs{z}.
\]

\section*{Exercício 6. [Intermediário --- produto próximo de 9]}

\textbf{Enunciado.} Se $\abs{x-3}<1/2$ e $\abs{y-3}<1/2$, prove que $\abs{xy-9}<13/4$.

\textbf{Solução.}
Expanda o desvio do produto:
\begin{align*}
(x-3)(y-3)+3(x-3)+3(y-3)
&=xy-3x-3y+9+3x-9+3y-9
=xy-9.
\end{align*}
A desigualdade triangular e a multiplicatividade do módulo dão
\begin{align*}
\abs{xy-9}
&\le\abs{x-3}\,\abs{y-3}+3\abs{x-3}+3\abs{y-3}.
\end{align*}
Falta estimar cada termo por uma desigualdade estrita. Sejam $p=\abs{x-3}$ e $q=\abs{y-3}$, de modo que $0\le p<1/2$ e $0\le q<1/2$.

Primeiro, $pq<(1/2)(1/2)=1/4$. Se $p=0$, então $pq=0<1/4$. Se $p>0$, de $q<1/2$ e $p>0$ segue $pq<p\cdot(1/2)$. De $p<1/2$ e $1/2>0$ segue $p\cdot(1/2)<(1/2)(1/2)$. A transitividade estrita dá $pq<1/4$.

Segundo, multiplicar $p<1/2$ e $q<1/2$ pelo positivo $3$ preserva o sentido: $3p<3/2$ e $3q<3/2$.

Some as três desigualdades estritas, um passo de cada vez. De $pq<1/4$ segue $pq+3p+3q<1/4+3p+3q$. De $3p<3/2$ segue $1/4+3p+3q<1/4+3/2+3q$. De $3q<3/2$ segue
\[
\frac14+\frac32+3q<\frac14+\frac32+\frac32=\frac14+3=\frac{13}{4}.
\]
Logo $\abs{xy-9}\le pq+3p+3q<13/4$.

\section*{Exercício 7. [Intermediário --- fator linear repetido]}

\textbf{Enunciado.} Decomponha $\dfrac{4x-1}{(x+1)^2}$ e verifique a identidade.

\textbf{Solução.}
Procura-se $A,B\in\R$ com
\[
\frac{4x-1}{(x+1)^2}=\frac{A}{x+1}+\frac{B}{(x+1)^2},\qquad x\neq -1.
\]
Multiplicando por $(x+1)^2$,
\[
4x-1=A(x+1)+B=Ax+(A+B).
\]
Igualando coeficientes,
\[
\begin{cases}
A=4,\\
A+B=-1.
\end{cases}
\]
A segunda equação dá $4+B=-1$, logo $B=-5$. A solução é única. Conferência:
\[
4(x+1)-5=4x+4-5=4x-1.
\]
Portanto
\[
\frac{4x-1}{(x+1)^2}=\frac{4}{x+1}-\frac{5}{(x+1)^2}.
\]

\section*{Exercício 8. [Avançado --- igualdade na triangular]}

\textbf{Enunciado.} Prove que $\abs{x+y}=\abs{x}+\abs{y}$ se, e somente se, $xy\ge 0$.

\textbf{Solução.}
Os dois lados de $\abs{x+y}=\abs{x}+\abs{y}$ são não negativos. Para números $p,q\ge 0$, tem-se $p=q$ se e somente se $p^2=q^2$: a ida é imediata da multiplicação; a volta usa $p^2-q^2=(p-q)(p+q)=0$, logo $p=q$ ou $p=-q$, e $p=-q$ com $p,q\ge 0$ força $p=q=0$.

Além disso $\abs{t}=t$ se e somente se $t\ge 0$, pela definição do módulo, e $\abs{xy}=\abs{x}\abs{y}$, e $\abs{x}^2=x^2$.

$(\Rightarrow)$ Suponha $\abs{x+y}=\abs{x}+\abs{y}$. Elevando ao quadrado,
\[
(x+y)^2=\bigl(\abs{x}+\abs{y}\bigr)^2,
\]
isto é
\[
x^2+2xy+y^2=\abs{x}^2+2\abs{x}\abs{y}+\abs{y}^2=x^2+2\abs{xy}+y^2.
\]
Subtraindo $x^2+y^2$, obtém-se $xy=\abs{xy}$. Pela caracterização acima, $xy\ge 0$.

$(\Leftarrow)$ Suponha $xy\ge 0$. Então $xy=\abs{xy}=\abs{x}\abs{y}$, e
\begin{align*}
(x+y)^2
&=x^2+2xy+y^2
=x^2+2\abs{x}\abs{y}+y^2
=\bigl(\abs{x}+\abs{y}\bigr)^2.
\end{align*}
Como $\abs{x+y}\ge 0$ e $\abs{x}+\abs{y}\ge 0$, a observação do primeiro parágrafo dá $\abs{x+y}=\abs{x}+\abs{y}$.

\section*{Exercício 9. [Avançado --- fator quadrático irredutível]}

\textbf{Enunciado.} Decomponha $\dfrac{3x^2+x+2}{(x-1)(x^2+x+1)}$ sobre $\R$, justificando a irredutibilidade e resolvendo o sistema.

\textbf{Solução.}
O discriminante de $x^2+x+1$ é $1-4=-3<0$. Um polinômio real de grau $2$ com discriminante negativo não tem raiz real e não se fatora em polinômios reais de grau $1$; é irredutível sobre $\R$.

Para $x\neq 1$, procura-se $A,B,C\in\R$ com
\[
\frac{3x^2+x+2}{(x-1)(x^2+x+1)}=\frac{A}{x-1}+\frac{Bx+C}{x^2+x+1}.
\]
Multiplicando pelo denominador,
\begin{align*}
3x^2+x+2
&=A(x^2+x+1)+(Bx+C)(x-1)\\
&=Ax^2+Ax+A+Bx^2-Bx+Cx-C\\
&=(A+B)x^2+(A-B+C)x+(A-C).
\end{align*}
O sistema é
\[
\begin{cases}
A+B=3,\\
A-B+C=1,\\
A-C=2.
\end{cases}
\]
A terceira equação equivale a $A=C+2$. Substituindo na primeira, $C+2+B=3$, logo $B=1-C$. Na segunda,
\begin{align*}
(C+2)-(1-C)+C
&=C+2-1+C+C
=3C+1.
\end{align*}
A equação $3C+1=1$ dá $3C=0$, logo $C=0$, $A=2$ e $B=1$. O coeficiente $3\neq 0$ torna o valor de $C$ único, e $A$ e $B$ ficam determinados. Conferência:
\[
2(x^2+x+1)+x(x-1)=2x^2+2x+2+x^2-x=3x^2+x+2.
\]
A decomposição é
\[
\frac{3x^2+x+2}{(x-1)(x^2+x+1)}=\frac{2}{x-1}+\frac{x}{x^2+x+1}.
\]

\end{document}
